\documentclass[10pt]{article}
\usepackage[margin=0.84in]{geometry}
\usepackage{amsmath,amssymb,mathtools}
\usepackage{booktabs,tabularx}
\usepackage[dvipsnames]{xcolor}
\usepackage[colorlinks=true,linkcolor=NavyBlue,urlcolor=BrickRed,hypertexnames=false]{hyperref}
\usepackage{microtype,fancyhdr,enumitem}
\definecolor{ink}{HTML}{3A293A}
\definecolor{accent}{HTML}{A44E82}
\definecolor{pale}{HTML}{FFF0F8}
\definecolor{passgreen}{HTML}{16794A}
\color{ink}
\pagestyle{fancy}\fancyhf{}
\lhead{$U(1)$ heat kernel}\rhead{Compact-group verification}\cfoot{\thepage}
\setlength{\headheight}{14pt}
\newcommand{\Sfun}{\mathcal S}
\newcommand{\PASS}{\textcolor{passgreen}{\textbf{PASS}}}
\newcommand{\summarybox}[1]{\begin{center}\fcolorbox{accent}{pale}{\parbox{0.91\linewidth}{#1}}\end{center}}
\title{\textbf{Heat-Kernel $\Lambda$-Distributions on $U(1)$}\\[3pt]
\large Casimir filtering, endpoint behavior, and alias-free quadrature}
\author{Proof and computational verification note}
\date{17 July 2026}

\begin{document}
\maketitle
\summarybox{\textbf{Outcome.}  Three diagonal $U(1)$ representations of
dimensions $1,2,3$ were checked at five heat times for all partitions through
total degree five.  Direct quadrature against the heat kernel agrees with the
weight/Casimir formula in all \textbf{285} cases.  The maximum absolute error
is $5.73\times10^{-14}$.}

\section{Heat measure and normalization}
Identify $U(1)$ with $e^{i\theta}$, $0\leq\theta<2\pi$, and choose the
Laplacian normalization for which the character $e^{iq\theta}$ has Casimir
$q^2$.  Brownian motion with generator $\tfrac12\Delta$ has density relative
to normalized Haar measure
\begin{equation}
k_t(\theta)=\sum_{q\in\mathbb Z}e^{-q^2t/2}e^{-iq\theta}.   \tag{1}
\end{equation}
Orthogonality immediately yields
\begin{equation}
\int_{U(1)}e^{iw\theta}k_t(\theta)\,d\theta/(2\pi)
=e^{-w^2t/2}.                                               \tag{2}
\end{equation}
This fixes the otherwise convention-dependent factor in the Casimir exponent.

Let $V$ be the diagonal representation with integer weights
$a_1,\ldots,a_d$,
\begin{equation}
\rho(e^{i\theta})=\operatorname{diag}
(e^{ia_1\theta},\ldots,e^{ia_d\theta}).                    \tag{3}
\end{equation}
For $W_\tau=\bigotimes_j\operatorname{Sym}^{\tau_j}V$, write
\begin{equation}
W_\tau=\bigoplus_{w\in\mathbb Z}M_\tau(w)\,\mathbb C_w.  \tag{4}
\end{equation}
Combining (2)--(4) proves the coefficient formula
\begin{equation}
\boxed{
[\mathbf t^\tau]\Sfun_t
=\sum_{w\in\mathbb Z}M_\tau(w)e^{-w^2t/2}.}               \tag{5}
\end{equation}
This is the compact-group heat formula in its simplest exact form: each
irreducible weight is attenuated by its Casimir eigenvalue.

\section{Endpoints and interpretation}
At $t=0$, every exponential in (5) is one, hence
\begin{equation}
[\mathbf t^\tau]\Sfun_0=\dim W_\tau,
\qquad
\Sfun_0(\mathbf t)=\prod_j(1-t_j)^{-d}.                    \tag{6}
\end{equation}
This is evaluation at the identity matrix.  As $t\to\infty$, only weight zero
survives:
\begin{equation}
\lim_{t\to\infty}[\mathbf t^\tau]\Sfun_t=M_\tau(0)
=\dim W_\tau^{U(1)}.                                       \tag{7}
\end{equation}
Thus (5) gives a literal spectral interpolation between dimension and Haar
invariants.  For a purely positive defining weight there are no positive-degree
Haar invariants.  Mixed positive and negative weights, however, can cancel and
leave nontrivial weight-zero terms.

\section{Representations and verification method}
The tested weight lists are
\begin{center}
\begin{tabular}{@{}lccp{7.1cm}@{}}
\toprule
Name & weights & dimension & purpose\\
\midrule
defining & $(1)$ & 1 & isolates the scalar Casimir $r^2$\\
weight pair & $(1,-1)$ & 2 & creates many weight-zero invariants\\
mixed & $(0,1,-2)$ & 3 & combines a fixed line with asymmetric weights\\
\bottomrule
\end{tabular}
\end{center}
For each representation, $t=0,0.25,0.8,3,20$ and all 19 partitions through
total degree five are tested.

The direct route uses 512 equally spaced angles.  At each angle it constructs
the matrix (3), computes complete characters from the matrix power traces, and
integrates their product against the Fourier heat kernel (1), truncated at
$|q|\leq32$.  The formula route never evaluates an angle or matrix: it builds
the integer multiplicities $M_\tau(w)$ and applies (5).

\subsection*{Why the quadrature is especially strong}
This is not a generic Monte Carlo comparison.  Through total degree five the
largest character frequency is at most $5\max_j|a_j|=10$.  Multiplication by
the truncated kernel produces frequencies at most 42.  Since 512 exceeds
twice this bandwidth, the periodic trapezoidal rule integrates the entire
trigonometric polynomial exactly in exact arithmetic.  At the smallest
positive time, the first omitted heat coefficient is
$e^{-33^2/8}<10^{-59}$.  The observed errors are therefore floating-point
roundoff, not sampling noise or appreciable heat-kernel truncation.

At $t=0$ the heat density is a delta distribution, so the code evaluates the
identity matrix directly as in (6).  This avoids approximating a singular
endpoint with a finite Fourier sum.

\section{Results}
\begin{center}
\begin{tabularx}{\linewidth}{@{}lXrrc@{}}
\toprule
Representation and time & coefficient & quadrature & Casimir formula & result\\
\midrule
defining, $t=0.8$ & $(3)$ & $0.0273237224$ & $0.0273237224$ & \PASS\\
weight pair, $t=3$ & $(2,1)$ & $0.8925233825$ & $0.8925233825$ & \PASS\\
mixed, $t=20$ & $(2,2)$ & $5.0003631994$ & $5.0003631994$ & \PASS\\
\bottomrule
\end{tabularx}
\end{center}
The last value is already close to the integer Haar multiplicity five, while
the small positive correction records the surviving nonzero weights at finite
time.

\section{Scope and reproducibility}
The computation verifies the compact heat formula for $U(1)$ with several
representation dimensions and both one-variable and mixed coefficients.  It
does not by itself verify the branching rules proposed for $U(n)$, $SU(n)$,
$O(n)$, or $USp(2n)$.  Those require separate highest-weight and Casimir
normalizations.  The disconnected-group warning is also essential: Brownian
motion from the identity in $O(n)$ remains in $SO(n)$ unless component-changing
jumps are added.

Run \path{compact_u1_heat/verification.py}, or open
\path{marimo_notebooks/u1_heat.py} in marimo.
\end{document}
